\clearpage
\section{판별식: 근의 충돌을 계수로 읽기}
\label{sec:polynomial-discriminant}

\begin{learninggoal}
반데르몬드 곱의 제곱을 이용하여 임의의 가환환 위 일계수다항식의 판별식을 정의한다. 근을 이용한 공식, 중근 판정과 교대군 판정을 증명한다.
\end{learninggoal}

변수 $T_1,\dots,T_n$에 대해
\[
  \delta(T_1,\dots,T_n)
  :=\prod_{1\le i<j\le n}(T_i-T_j)
\]
를 \emph{반데르몬드 곱}(Vandermonde product)이라 한다. 순열 $\pi\in S_n$은
\[
  \pi(\delta)=\sgn(\pi)\delta
\]
로 작용한다. 따라서 제곱
\[
  \delta^2=\prod_{i<j}(T_i-T_j)^2
\]
은 대칭다항식이다.

기본대칭식 정리에 의해 유일한
\[
  D_n(S_1,\dots,S_n)\in\Z[S_1,\dots,S_n]
\]
가 존재하여
\[
  \delta^2=D_n(s_1,\dots,s_n)
\]
이다. 이 보편다항식이 판별식을 계수만으로 정의하게 해 준다.

\begin{definition}[일계수다항식의 판별식]
가환환 $R$ 위의 일계수 $n$차다항식
\[
  f(X)=X^n+c_1X^{n-1}+\cdots+c_n
\]
의 \emph{판별식}(discriminant)을
\[
  \Disc(f)
  :=D_n(-c_1,c_2,\dots,(-1)^nc_n)
  \in R
\]
로 정의한다. $n=0$인 $f=1$에 대해서는 빈 곱의 관례에 따라 $\Disc(1)=1$로 둔다.
\end{definition}

이 정의에서 일반 계수환 위의 기본대칭식 정리가 필요하다. 먼저 정수계수 보편다항식 $D_n$을 만든 뒤, 환준동형
\[
  \Z[S_1,\dots,S_n]\longrightarrow R,
  \qquad
  S_j\longmapsto(-1)^jc_j
\]
으로 옮기기 때문이다.

\begin{theorem}[근을 이용한 판별식 공식]
\label{thm:discriminant-root-formula}
$R\subseteq R'$가 가환환의 확장이고
\[
  f(X)=\prod_{i=1}^n(X-\alpha_i)
  \qquad(\alpha_i\in R')
\]
라 하자. 그러면
\[
  \Disc(f)
  =\prod_{1\le i<j\le n}(\alpha_i-\alpha_j)^2.
\]
근들은 중복도까지 포함하여 나열한다.
\end{theorem}

\begin{proof}
대입준동형
\[
  \Z[T_1,\dots,T_n]\longrightarrow R',
  \qquad
  T_i\longmapsto\alpha_i
\]
을 생각한다. 기본대칭식 $s_j$는 비에타 공식에 의해 $(-1)^jc_j$로 보내지고, $D_n(s_1,\dots,s_n)=\delta^2$이므로 원하는 등식이 나온다.
\end{proof}

\begin{corollary}[중근 판정]
$K$가 체이고 $f\in K[X]$가 일계수다항식이면 다음은 동치이다.
\begin{enumerate}[label=(\roman*)]
  \item $\Disc(f)=0$이다.
  \item $f$는 대수적 폐포에서 중근을 갖는다.
  \item $f$와 형식적 도함수 $f'$가 공통근을 갖는다.
  \item $\gcd(f,f')\ne1$이다.
\end{enumerate}
따라서
\[
  f\text{가 분리다항식}
  \quad\Longleftrightarrow\quad
  \Disc(f)\ne0.
\]
\end{corollary}

\begin{proof}
대수적 폐포에서 근들을 $\alpha_1,\dots,\alpha_n$이라 하면 근 공식에 의해 판별식은 근 차이들의 제곱곱이다. 체에서는 곱이 $0$인 것은 어떤 인수가 $0$인 것과 동치이므로, 판별식이 사라지는 것은 두 근이 같은 것과 같다. 나머지 동치는 중근과 도함수의 표준 판정에서 따른다.
\end{proof}

\begin{example}[이차다항식]
\[
  f=X^2+aX+b
\]
의 근을 $\alpha,\beta$라 하면
\[
  \Disc(f)=(\alpha-\beta)^2
  =(\alpha+\beta)^2-4\alpha\beta
  =a^2-4b.
\]
따라서 표수가 $2$가 아닌 체에서 $f$가 중근을 갖는 것은 $a^2-4b=0$인 것과 동치이다.
\end{example}

\subsection{판별식의 제곱근과 짝순열}

$f\in K[X]$가 분리인 $n$차다항식이고 $L$이 분해체라 하자. 근들을 $\alpha_1,\dots,\alpha_n$이라 쓰고
\[
  \delta_f:=\prod_{i<j}(\alpha_i-\alpha_j)\in L
\]
라 두면
\[
  \delta_f^2=\Disc(f).
\]
근의 번호를 바꾸면 $\delta_f$의 부호가 바뀔 수 있지만, 제곱은 변하지 않는다.

\begin{theorem}[교대군 판정]
\label{thm:discriminant-alternating-criterion}
$\charac K\ne2$이고 $f\in K[X]$가 분리인 $n$차다항식이라 하자. 근의 순열표현으로
\[
  G=\Gal(f/K)\le S_n
\]
라 보자. 그러면 다음은 동치이다.
\begin{enumerate}[label=(\roman*)]
  \item $G\subseteq A_n$이다.
  \item $\delta_f\in K$이다.
  \item $\Disc(f)$가 $K$에서 제곱이다.
\end{enumerate}
\end{theorem}

\begin{proof}
$\sigma\in G$가 근들에 유도하는 순열을 같은 기호로 쓰면
\[
  \sigma(\delta_f)=\sgn(\sigma)\delta_f.
\]
따라서 모든 $\sigma$가 짝순열인 것은 모든 $\sigma$가 $\delta_f$를 고정하는 것과 동치이고, 유한 갈루아 확장에서 전체 갈루아군의 고정체는 $K$이므로 (i)와 (ii)가 동치이다.

(ii)$\Rightarrow$(iii)는 자명하다. 반대로 $d^2=\Disc(f)$인 $d\in K$가 있다고 하자. $\delta_f$와 $d$는 $L$에서 $X^2-\Disc(f)$의 근이다. 표수가 $2$가 아니므로 두 근은 $\pm d$뿐이고, 따라서 $\delta_f=\pm d\in K$이다.
\end{proof}

\begin{remark}
삼차 기약다항식에서는 추이적 부분군이 $A_3\cong C_3$ 또는 $S_3$뿐이므로, 위 정리가 갈루아군을 완전히 결정한다. 더 높은 차수에서는 판별식은 갈루아군이 $A_n$ 안에 들어가는지만 알려 주며, 그 안에서 어느 부분군인지는 추가 정보가 필요하다.
\end{remark}

\subsection{변수의 평행이동과 확대}

\begin{proposition}
$f\in R[X]$가 일계수 $n$차다항식이고 $c\in R$라 하자. 그러면
\[
  g(X)=f(X+c)
\]
도 일계수이고
\[
  \Disc(g)=\Disc(f).
\]
또한 $u\in R^\times$가 단위원이고
\[
  h(X)=u^{-n}f(uX)
\]
라 하면
\[
  \Disc(h)=u^{-n(n-1)}\Disc(f).
\]
\end{proposition}

\begin{proof}
확장환에서 $f$의 근이 $\alpha_i$이면 $g$의 근은 $\alpha_i-c$이다. 근 차이는 변하지 않으므로 첫 등식이 따른다. $h$의 근은 $u^{-1}\alpha_i$이고 근 차이마다 $u^{-1}$이 곱해진다. 쌍의 수가 $n(n-1)/2$이므로 제곱곱에는 $u^{-n(n-1)}$이 나타난다.
\end{proof}

\begin{pitfall}
판별식은 근들의 곱이나 합이 아니라 \emph{서로 다른 두 근 사이의 차이}를 측정한다. 판별식이 $0$이라는 것은 근 하나가 $0$이라는 뜻이 아니라, 두 근이 충돌했다는 뜻이다.
\end{pitfall}

\begin{checkpoint}
\begin{enumerate}[label=(\alph*)]
  \item $X^3-3X+2$의 판별식이 $0$임을 근의 중복도를 이용하여 확인하라.
  \item $\Disc(f)$가 근의 번호 붙이기에 의존하지 않는 이유를 설명하라.
  \item 표수가 $2$일 때 교대군 판정의 증명이 어디에서 무너지는지 찾으라.
\end{enumerate}
\end{checkpoint}
